% !TeX program = lualatex
% !TeX encoding = utf8
% !TeX spellcheck = uk_UA
% !BIB program = biber

\documentclass[%
biblatex,
%marginversioninfo
]{ConspectBook}


\begin{document}

\subsection*{Гіперболічний поворот}
Розглянемо ще раз матрицю одновимірного перетворення Лоренца (буста):

\[
\|L^{\mu'}{}_{\mu}\| =
\begin{pmatrix}
\gamma & -\beta\gamma & 0 & 0 \\
-\beta\gamma & \gamma & 0 & 0 \\
0 & 0 & 1 & 0 \\
0 & 0 & 0 & 1
\end{pmatrix}, \qquad
\beta = \frac{v}{c},\quad \gamma = \frac{1}{\sqrt{1-\beta^2}}.
\]

Незважаючи на зовнішню схожість із матрицею просторового повороту в площині $(x^1, x^2)$,
\[
R(\theta) =
\begin{pmatrix}
1 & 0 & 0 & 0 \\
0 & \cos\theta & -\sin\theta & 0 \\
0 & \sin\theta & \cos\theta & 0 \\
0 & 0 & 0 & 1
\end{pmatrix},
\]
буст має принципову відмінність --- він зберігає не суму квадратів координат, а їх різницю.
Цю відмінність можна зробити наочною, ввівши замість швидкості $v$ \textbf{швидкісний параметр}
(rapidity) $\psi$:
\[
\beta = \th\psi, \qquad \gamma = \ch\psi, \qquad \beta\gamma = \sh\psi.
\]
Тоді матриця буста набуває вигляду
\begin{equation}
\|L^{\mu'}{}_{\mu}\| =
\begin{pmatrix}
\ch\psi & -\sh\psi & 0 & 0 \\
-\sh\psi &  \ch\psi & 0 & 0 \\
0 & 0 & 1 & 0 \\
0 & 0 & 0 & 1
\end{pmatrix}.
\label{eq:boost_hyperbolic}
\end{equation}
Це --- \textbf{гіперболічний поворот} у площині $(x^0, x^1)$, який залишає інваріантною
гіперболу $(x^0)^2 - (x^1)^2 = \const$, подібно до того, як звичайний поворот залишає
інваріантним коло. Ізотропія простору гарантує, що залежність $\psi(v)$ є непарною
($\psi(-v) = -\psi(v)$), а з властивостей гіперболічних функцій випливає
\textbf{адитивність} швидкісного параметра при послідовних бустах:
\[
\psi_{13} = \psi_{12} + \psi_{23},
\]
що еквівалентно формулі додавання швидкостей (див. підрозділ про додавання швидкостей).
У цьому сенсі $\psi$ є природнішою мірою руху, ніж $v$, оскільки вона пробігає всі дійсні
значення $(-\infty, +\infty)$, тоді як $v\in(-c,c)$.

Отже, перетворення Лоренца слід розуміти як \textbf{псевдоповороти} в просторі Мінковського,
які разом із звичайними просторовими обертаннями утворюють групу Лоренца. Будь-яке власне
перетворення Лоренца можна подати у вигляді композиції просторового обертання, буста вздовж
осі $x^1$ та ще одного просторового обертання.


\begin{figure}[h!]\centering
	\begin{tikzpicture}[
        declare function={atanh(\x) = 1/2*ln((1 + \x)/(1 -\x));}
    ]
		\pgfmathsetmacro{\hangle}{atanh(0.2)}
		\pgfmathsetmacro{\angle}{\hangle*45}

		\pgfmathdeclarefunction{x}{2}{\pgfmathparse{#1*cosh(\hangle)+#2*sinh(\hangle)}} % Real x coord of canvas
		\pgfmathdeclarefunction{y}{2}{\pgfmathparse{#1*sinh(\hangle)+#2*cosh(\hangle)}} % Real y coord of canvas
		\def\point(#1,#2){({x(#1,#2)},{y(#1,#2)})}

		\draw [-latex',thick, black] (0,0) -- (10,0)  node[below] {$x$};
		\draw [-latex',thick, black] (0,0) -- (0,10)  node[left] {$t$};
		\foreach \i in {1,...,10} {
			\draw [dashed, black] (0,\i) -- (10,\i)  ;
			\draw [dashed, black] (\i,0) -- (\i,10)  ;
		}

		\draw [-latex',thick, red] \point(0,0) -- \point(10,0)  node[below] {$x'$};
		\draw [-latex',thick, red] \point(0,0) -- \point(0,10)  node[left] {$t'$};
		\foreach \i in {1,...,10} {
			\draw [dashed, red!50] \point(0,\i) -- \point(10,\i)  ;
			\draw [dashed, red!50] \point(\i,0) -- \point(\i,10)  ;
		}

%		\fill [gray] (3,6) circle (0.1) node [above] {$\mathbf{A}$}
%		(5,6) circle (0.1) node [above] {$\mathbf{B}$};
%
%		\fill[red] \point(3,6) circle (0.1) node [above] {$\mathbf{A}'$}
%		\point(5,6) circle (0.1) node [above] {$\mathbf{B}'$};
	\end{tikzpicture}
	\caption{}
\end{figure}

\end{document}
